\begin{solucion}
\subsection*{Lado izquierdo}
Sea
\[
  F[n_,k_] := \bigl(\binom{n}{k}\bigr)^{3}.
\]

El siguiente bloque crea la función y comprueba que la suma es \emph{hipergeométrica} en $k$.

\begin{lstlisting}[caption={Definición de \texttt{F} y aplicación de Gosper}]
In[12]:= F[n_, k_] := (Binomial[n, k])^3
In[13]:= Gosper[F[n, k], k]
Out[13]= {}
\end{lstlisting}

Una salida vacía de \textsc{Gosper} indica que no existe una antiderivada hipergeométrica—exactamente el escenario
que permite a \textsc{Zeilberger} generar una recurrencia para la suma.

\begin{lstlisting}[caption={Recurrencia obtenida vía \textsc{Zeilberger}}]
In[14]:= Zb[F[n, k], k, n]
Out[14]= {-8 (1 + n)^2 f[n] + (-16 - 21 n - 7 n^2) f[1 + n] + (2 + n)^2 f[2 + n] == 0}
\end{lstlisting}

Guardamos la ecuación como \texttt{recurrenciaF} y pedimos a \texttt{Hyper} verificar que la sucesión
$S_F(n)=\sum_{k=0}^{n} F[n,k]$ la satisface.

\begin{lstlisting}[caption={Verificación automática de la recurrencia}]
In[15]:= recurrenciaF = -8 (1 + n)^2 f[n] + (-16 - 21 n - 7 n^2) f[1 + n] + (2 + n)^2 f[2 + n] == 0;
In[16]:= Hyper[recurrenciaF, f[n]]
Out[16]= {}
\end{lstlisting}

\subsection*{Lado derecho}
Definimos ahora
\[
  H[n_,k_] := \binom{n}{k}^{2}\,\binom{2k}{n}.
\]

\begin{lstlisting}[caption={Definición de \texttt{H} y aplicación de Gosper}]
In[17]:= H[n_, k_] := (Binomial[n, k])^2*Binomial[2 k, n]
In[18]:= Gosper[H[n, k], k]
Out[18]= {}
\end{lstlisting}

Aplicando \textsc{Zeilberger}:

\begin{lstlisting}[caption={Recurrencia para el lado derecho}]
In[19]:= Zb[H[n, k], k, n]
Out[19]= {-8 (1 + n)^2 f[n] + (-16 - 21 n - 7 n^2) f[1 + n] + (2 + n)^2 f[2 + n] == 0}
\end{lstlisting}

Observamos que la \emph{misma} recurrencia aparece para $S_H(n)=\sum_{k=0}^{n} H[n,k]$.  Al igual que antes:

\begin{lstlisting}
In[20]:= Hyper[recurrenciaF, f[n]]
Out[20]= {}
\end{lstlisting}

% -----------------------------------------------------------------------------
\section{Coincidencia de valores iniciales}

Para concluir que $S_F(n)=S_H(n)$ para todo $n\ge 0$ necesitamos comprobar que ambas sucesiones coinciden en
los dos primeros valores, pues la recurrencia es de orden $2$.

\begin{center}
\begin{tabular}{c|c|c}
$n$ & $S_F(n)$ & $S_H(n)$ \\ \hline
0 & 1 & 1 \\
1 & 2 & 2 \\
\end{tabular}
\end{center}

Los cómputos son directos y se muestran a continuación.

\begin{lstlisting}
In[21]:= Table[{n, Sum[F[n, k], {k, 0, n}], Sum[H[n, k], {k, 0, n}]}, {n, 0, 1}]
Out[21]= {{0, 1, 1}, {1, 2, 2}}
\end{lstlisting}

% -----------------------------------------------------------------------------
\section{Conclusión}

Puesto que $S_F$ y $S_H$ satisfacen la misma ecuación lineal homogénea de segundo orden y concuerdan en sus
valores iniciales, se sigue por el principio de inducción que $S_F(n)=S_H(n)$ para todo $n\ge 0$.  Por
\eqref{eq:main}, la identidad queda demostrada.

\bigskip
\noindent\textbf{Archivo de respaldo.} El cuaderno \texttt{Taller7\_8.nb} con todo el código ejecutado se incluye
como material complementario.
\end{solucion}